\chapter*{Resumo}
\addcontentsline{toc}{chapter}{Resumo}
\markboth{Resumo}{}

A produção de biodiesel por transesterificação alcalina homogênea, rota
predominante na indústria, esbarra em dificuldades conhecidas de separação e
purificação, além de gerar efluentes que precisam ser neutralizados. A
substituição do catalisador dissolvido por um sólido elimina boa parte desses
inconvenientes, mas introduz em seu lugar um problema de transporte: os
reagentes passam a depender da difusão até a superfície ativa, e o óleo vegetal,
cuja difusividade mássica é cinco ordens de grandeza menor que a difusividade de
momento, não é um meio favorável a esse tipo de transporte.

Reatores monolíticos de parede catalítica oferecem uma resposta possível a esse
problema, por combinarem elevada razão entre área superficial e volume com perda
de carga reduzida. A pergunta que o presente trabalho procura responder é se a
geometria monolítica de fato se justifica para a transesterificação do óleo de
soja com óxido de cálcio, ou se o sistema opera em um regime no qual a área
adicional não é aproveitada.

O trabalho é conduzido por meio de simulação fluidodinâmica computacional de um
canal representativo do monólito, em regime laminar e isotérmico a \SI{65}{\celsius},
com o catalisador tratado como condição de contorno de fluxo de espécies na
parede. A resistência difusional no interior do \textit{washcoat} é incorporada
por um fator de efetividade calculado a partir do módulo de Thiele, o que evita
discretizar a camada porosa. Os parâmetros cinéticos provêm de medidas
realizadas no próprio laboratório com CaO suportado em biochar.

Os resultados preliminares indicam que o número de Damköhler de parede vale
\num{0.76}, o que situa o sistema em regime misto: a resistência cinética e a
resistência ao transporte externo são da mesma ordem, e a concentração de
triglicerídeo junto à superfície catalítica é \SI{43}{\percent} menor que no seio
do fluido. Esse valor foi reproduzido em duas malhas independentes e em duas
composições de alimentação distintas, com discrepância inferior a
\SI{2}{\percent}. Verificou-se ainda que o comprimento de entrada mássica do
canal simulado é de \SI{3.46}{\meter}, contra os \SI{50}{\milli\meter} de
comprimento físico, de modo que a camada-limite de concentração jamais se
desenvolve e a conversão por passe permanece baixa, em torno de
\SI{1.4}{\percent}.

Nas etapas seguintes, o modelo será calibrado contra os dados experimentais em
obtenção no laboratório, estendido à varredura de temperatura e de espessura de
\textit{washcoat}, e empregado para a avaliação de geometrias alternativas de
canal.

\vspace{10pt}
\noindent\textbf{Palavras-chave:} biodiesel; transesterificação heterogênea;
reator monolítico; fluidodinâmica computacional; número de Damköhler.
